\subsectionaddtolist{۱.۷ تحلیل دامنه نفوذ}

{آ)} خیر، مهاجم نمی‌تواند یک گواهی معتبر برای یک سرور جعلی \lr{EHR} جعل کند. سرورهای \lr{EHR} زیر نظر \lr{Infrastructure Intermediate CA} هستند، در حالی که کلید خصوصی فاش‌شده مربوط به \lr{User \& Device Intermediate CA} است. از آنجایی که این دو مرجع صدور مجزا هستند، کلید فاش‌شده توانایی امضا و صدور گواهی برای زیرساخت شبکه و سرورهای \lr{EHR} را ندارد.

مهاجم تنها می‌تواند گواهی‌های جعلی در دامنه عملکرد \lr{User \& Device Intermediate CA} صادر کند که شامل موارد زیر است:
\begin{itemize}
    \item گواهی‌های هوشمند شخصی کارکنان (پزشکان و پرستاران) برای ورود به سیستم‌ها یا امضای دیجیتال ایمیل‌ها (\lr{S/MIME}).
    \item گواهی‌های مربوط به دستگاه‌های اینترنت اشیاء پزشکی (\lr{IoT}).
\end{itemize}

{ب)} ارزیابی ویژگی‌های امنیتی در حوزه عملکرد کاربران و دستگاه‌های \lr{IoT}:
\begin{itemize}
    \item {صحت (\lr{Integrity}):} مهاجم با در دست داشتن کلید خصوصی این \lr{CA} می‌تواند گواهی‌های جعلی \lr{S/MIME} صادر کند و ایمیل‌های حاوی اطلاعات بیماران را به نام پزشکان یا کارکنان مجاز امضا کند. در نتیجه، گیرنده پیام گمان می‌کند که محتوا دست‌نخورده و معتبر است، در حالی که محتوای ایمیل توسط مهاجم تغییر یافته یا جعل شده است.
    \item {محرمانگی (\lr{Confidentiality}):} مهاجم می‌تواند با صدور گواهی‌های جعلی برای دستگاه‌های \lr{IoT} پزشکی یا سیستم‌های کاربران، خود را به عنوان یک گیرنده مجاز جا بزند. به عنوان مثال، اگر کاربران ایمیل‌های رمزگذاری‌شده با \lr{S/MIME} ارسال کنند، مهاجم با گواهی جعلی مرتبط با آن کلید عمومی، قادر به رمزگشایی و شنود اطلاعات حساس بیماران خواهد بود.
    \item {احراز هویت (\lr{Authentication}):} گواهی‌های احراز هویت روی کارت‌های هوشمند کارکنان و گواهی‌های تعبیه‌شده در سخت‌افزار دستگاه‌های \lr{IoT} برای احراز هویت متقابل با سرورها بر بستر \lr{TLS} استفاده می‌شوند. مهاجم با جعل این گواهی‌ها می‌تواند به عنوان یک پزشک، پرستار یا یک دستگاه \lr{IoT} معتبر به سیستم‌های بیمارستان متصل شده و فرآیند احراز هویت را به طور کامل دور بزند.
    \item {عدم انکار (\lr{Non-repudiation}):} از آنجایی که امضای دیجیتال ایمیل‌ها بر پایه این گواهی‌ها بنا شده است، اگر مهاجم ایمیلی مخرب یا دستور درمانی اشتباهی را با یک گواهی جعلی امضا و ارسال کند، سیستم آن را به عنوان امضای معتبر پزشک می‌شناسد. این امر باعث می‌شود که پزشک واقعی نتواند ارسال آن پیام را انکار کند و مسئولیت اقدامات مهاجم به اشتباه متوجه کارکنان شود.
\end{itemize}

{ج)} خیر، نیازی به باطل کردن گواهی \lr{Root CA} یا گواهی‌های صادر شده توسط \lr{Infrastructure Intermediate CA} وجود ندارد. در ساختار سلسله‌مراتب مراجع صدور (\lr{CA Hierarchy})، شاخه‌ها مستقل از یکدیگر عمل می‌کنند. از آنجا که نفوذ تنها به سطح کلید خصوصی \lr{User \& Device Intermediate CA} محدود بوده و به کلیدهای بالا دست (\lr{Root CA}) یا شاخه‌های موازی سرایت نکرده است، گواهی‌های بخش زیرساخت و سرورهای شبکه همچنان امن و معتبر باقی می‌مانند. تنها باید گواهی خودِ \lr{User \& Device Intermediate CA} توسط \lr{Root CA} باطل شود.


\subsectionaddtolist{۲.۷ ابطال گواهی‌های جعل‌شده}

{آ)} چالش‌ها و محدودیت‌های عملیاتی استفاده از لیست ابطال گواهی (\lr{CRL}):
\begin{itemize}
    \item {حجم فایل:} با توجه به تعداد بسیار زیاد دستگاه‌های \lr{IoT} و کلاینت‌های کاربری، باطل کردن انبوه گواهی‌ها باعث می‌شود حجم فایل \lr{CRL} به شدت افزایش یابد.
    \item {پهنای باند شبکه:} دانلود دوره‌ای یک فایل \lr{CRL} حجیم توسط هزاران دستگاه \lr{IoT} و سیستم کاربری به صورت همزمان، پهنای باند شبکه بیمارستان را اشباع کرده و کارایی ارتباطات حیاتی را کاهش می‌دهد.
    \item {تاخیر (\lr{Latency}):} دستگاه‌ها باید پیش از هر ارتباط، کل فایل \lr{CRL} را دانلود و در میان لیست طولانی آن جستجو کنند که این امر تاخیر (\lr{Latency}) فرآیند احراز هویت را به شدت افزایش می‌دهد. همچنین فرآیند به‌روزرسانی \lr{CRL} زمان‌بر است و در پنجره زمانی انتشار آن، گواهی‌های جعل‌شده ممکن است همچنان معتبر شناخته شوند.
\end{itemize}

{ب)} پروتکل \lr{OCSP} به جای ارسال کل لیست ابطال، به کلاینت‌ها اجازه می‌دهد وضعیت یک گواهی خاص را به صورت لحظه‌ای و در قالب یک پاسخ کوچک از سرور استعلام کنند. این کار چالش حجم فایل و پهنای باند را در مقایسه با \lr{CRL} پوشش می‌دهد.

{نقاط ضعف احتمالی \lr{OCSP}:}
\begin{itemize}
    \item {نقطه شکست واحد (\lr{SPOF}) و بار پردازشی سنگین:} مواجهه با حجم عظیم درخواست‌های همزمان از سوی تجهیزات پزشکی، سرور \lr{OCSP Responder} را زیر بار شدید می‌برد و ممکن است منجر به حملات ناخواسته محرومیت از سرویس (\lr{DoS}) شود.
    \item {افزایش تاخیر خط ارتباطی:} برای هر اتصال \lr{TLS}، دستگاه باید یک استعلام آنلاین اضافه به سرور \lr{OCSP} بفرستد و منتظر پاسخ بماند که در تجهیزات پزشکی حساس زمان‌بر است.
    \item {حریم خصوصی:} سرور \lr{OCSP} متوجه می‌شود که کدام دستگاه در چه زمانی به چه سرویسی متصل شده است.
\end{itemize}
\textit{راهکار بهبود:} برای رفع این نقاط ضعف می‌توان از \lr{OCSP Stapling} استفاده کرد تا خودِ سرورها وضعیت ابطال گواهی‌شان را ضمیمه کنند و بار از روی سرور استعلام برداشته شود.


\subsectionaddtolist{۳.۷ بهبودهای پیشنهادی معماری \lr{PKI}}

\begin{enumerate}
    \item {استفاده از مکانیزم \lr{OCSP Stapling}:} برای کاهش بار سرورهای پاسخ‌دهنده و از بین بردن تاخیر استعلام آنلاین دستگاه‌های \lr{IoT}، وضعیت اعتبار گواهی به همراه پاسخ دست‌تکانی \lr{TLS} توسط خودِ سرور به کلاینت تحویل داده شود.
    \item {جداسازی دامنه‌های امنیتی به مراجع صدور میانی بیشتر (\lr{Compartmentalization}):} دستگاه‌های \lr{IoT} با ریسک بالا را از بخش کاربران عادی تفکیک کرده و برای هر کدام یک \lr{Sub-CA} مجزا در نظر گرفته شود تا در صورت لو رفتن یک کلید، دامنه نفوذ بسیار کوچک‌تر باشد.
    \item {به‌کارگیری ماژول‌های سخت‌افزاری امنیت (\lr{HSM}):} کلیدهای خصوصی تمامی مراجع صدور میانی (\lr{Intermediate CAs}) باید به جای ذخیره‌سازی روی سیستم‌عامل سرور، درون سخت‌افزارهای اختصاصی ضد دستکاری (\lr{HSM}) نگهداری شوند تا امکان استخراج و فاش شدن کلیدها توسط باج‌افزارها به صفر برسد.
\end{enumerate}